% ============================================================================
% Fichier  : partie5/P05_C15_risques_emergents.tex
% Titre    : Risques Émergents et Nouveaux Objets d'Audit
% Auteur   : MINKA MI NGUIDJOÏ Thierry Emmanuel
% Version  : 1.0 -- novembre 2026
% Corrigés : corriges/Corr_C15.tex
% ============================================================================
% Sources : ISO/IEC 27002:2022 (mesures 5.19 à 5.23, intitulés traduits) ;
% ISO/IEC 27017:2015 et 27018:2019 ; ISO/IEC 42001:2023 (système de
% management de l'IA, certifiable) ; NIST AI RMF 1.0 (2023 : fonctions
% Gouverner, Cartographier, Mesurer, Gérer ; caractéristiques de
% confiance) ; ISAE 3402 / SOC 2 (Chapitre 9). Loi n° 2010/012, art. 26
% al. 1 (cité au Chapitre 3) ; loi n° 2024/017 à vérifier.
% Faits du corpus : ANPP (quatre agents payés après leur départ) ; LiZbiZ
% (148 M FCFA d'ajustements de stock, dont 71 % en fin de mois, sous le
% compte de M. FOUDA ; cours ACSI).
% Faits nouveaux (reportés dans la bible) : en 2029, la SANAGA migre sa
% messagerie vers un service en nuage (NUAGEPRO, fictif) ; registre des
% 9 prestataires critiques (4/1/4) ; modèle de détection de fraude aux
% virements acheté à un éditeur (FRAUDSCAN, fictif) ; usage d'un assistant
% d'IA générative public par des agents (copie de données de clients) ;
% audit continu des paiements de l'ANPP : 37 exceptions en juin 2029
% (29 justifiées, 6 indues, 2 inconnues).
% ============================================================================

\chapter{Risques Émergents et Nouveaux Objets d'Audit}
\label{chap:emergents}
\chaptermark{Risques émergents}

\epigraph{%
    « On peut externaliser un service.\\
    On n'externalise jamais la responsabilité. »%
}{Principe du cours}

\niveauChapitre{Expert}{%
    Maîtrise du Chapitre~\ref{chap:cas-integrateurs}%
}

\begin{boxobjectifs}
\begin{itemize}
    \item Gérer et auditer les risques liés aux tiers et à la chaîne
          d'approvisionnement.
    \item Auditer un service en nuage selon le partage des responsabilités.
    \item Situer les risques des systèmes d'intelligence artificielle, et les
          cadres qui les encadrent (ISO/IEC~42001, NIST AI RMF).
    \item Concevoir un dispositif d'audit continu, et en connaître les
          limites.
    \item Appliquer à ces nouveaux objets la cote \OUI{}~/ \NON{}~/ \NEK{} et
          les trois mesures.
\end{itemize}
\end{boxobjectifs}

\bigskip

% ============================================================================
\section*{Avant de commencer : la sécurité qui n'est plus chez soi}
% ============================================================================

Douala, 2029. La BANQUE SANAGA\footnote{Comme dans tout le cours, les
organisations, les prestataires (DIGICORE, NUAGEPRO, FRAUDSCAN), les
personnes, les dates et les montants sont fictifs~; les normes et les
cadres cités sont exacts.} a reconstruit sa sécurité sur ce qu'elle
détient~: ses postes, ses réseaux, ses agences. Mais une part croissante de
son système d'information n'est plus chez elle. Le centre de données est
hébergé par DIGICORE~; la messagerie migre en 2029 vers un service en nuage,
NUAGEPRO~; un modèle d'intelligence artificielle, acheté à l'éditeur
FRAUDSCAN, analyse désormais chaque virement. Et des agents ont pris
l'habitude de coller des extraits de dossiers de clients dans un assistant
d'IA générative public, «~pour rédiger plus vite~».

La RSSI résume la situation au comité~: «~Nous savons auditer nos systèmes.
Nous ne savons pas encore auditer ceux des autres, ni ceux qui
apprennent.~»

\medskip

Ce chapitre ouvre la Partie~V. Il présente trois objets d'audit qui
transforment le métier~: les tiers et le nuage, l'intelligence
artificielle, et l'audit lui-même, qui devient continu.

\begin{boxalert}
\textbf{Violation typique.} Considérer qu'un risque confié à un prestataire
sort du registre de l'organisation.

\smallskip
\textbf{Règle~:} le risque reste celui de l'organisation~; seule
l'exécution de certaines mesures change de main. Le registre porte le
risque, le contrat porte les obligations du prestataire, et la preuve doit
revenir à l'organisation.
\end{boxalert}

\separateur

% ============================================================================
\section{Les Tiers et la Chaîne d'Approvisionnement}
\label{sec:c15-tiers}
% ============================================================================

\subsection{Un risque qui vient de l'extérieur}

Le Chapitre~\ref{chap:ebios} a montré qu'un attaquant ciblé passe souvent
par un partenaire moins protégé. Les tiers posent trois questions
distinctes~: un prestataire peut \emph{servir de porte d'entrée} (l'accès de
maintenance de LiZbiZ)~; il peut \emph{défaillir} (DIGICORE)~; et il peut
lui-même dépendre d'autres prestataires, que l'organisation ne connaît pas
(les sous-traitants de rang deux).

L'ISO/IEC~27002:2022 consacre cinq mesures à ces sujets~: la politique de
sécurité dans les relations avec les fournisseurs (5.19), la sécurité dans
les accords (5.20), la chaîne d'approvisionnement des TIC (5.21), la
surveillance et la revue des services des fournisseurs (5.22), et la
sécurité des services en nuage (5.23)~\cite{iso27002_2022}.

\subsection{Le cycle de vie d'un tiers}

\begin{tableaucours}{L{2.6cm}L{9.6cm}}{%
    \tableauHeader{Étape} &
    \tableauHeader{Ce que l'organisation fait, et ce qu'elle garde comme preuve}
}
Sélection & Classe le tiers selon sa criticité (dépendance, pénétration~:
Chapitre~\ref{chap:ebios})~; évalue sa maturité par un questionnaire et des
attestations \\
Contrat & Inscrit les exigences de sécurité, la notification des
incidents, le droit d'audit ou le rapport d'assurance, la réversibilité,
l'accord préalable sur la sous-traitance \\
Exploitation & Suit les indicateurs de service et de sécurité, reçoit les
rapports d'assurance, exerce son droit d'audit~; met à jour la criticité \\
Sortie & Récupère les données, révoque les accès, obtient la preuve de
destruction \\
\end{tableaucours}
\vspace{-0.8em}
\begin{center}{\small\itshape Tableau~15.1. Le cycle de vie d'un tiers}\end{center}

\begin{boxlecture}{Les neuf prestataires critiques de la BANQUE SANAGA}
\textbf{Constatation.} En mars 2029, M.~ATANGANA cote, pour chacun des neuf
prestataires classés critiques, l'existence d'une preuve récente que ses
mesures fonctionnent (rapport d'assurance de type~2, audit exercé, ou
certification couvrant le service)~:
\begin{itemize}
    \item \OUI{}~: le prestataire des GAB (audit exercé grâce à
          l'avenant), le fournisseur d'accès (test de la protection contre
          le déni de service), l'éditeur du système central (rapport de
          type~2), NUAGEPRO (certification ISO/IEC~27001 et rapport de
          type~2 couvrant la messagerie)~;
    \item \NON{}~: DIGICORE (rapport annuel non remis)~;
    \item \NEK{}~: FRAUDSCAN, le prestataire de la monétique, l'imprimeur
          des chéquiers et le transporteur de fonds (aucune preuve demandée
          à ce jour).
\end{itemize}
4~\OUI{}, 1~\NON{}, 4~\NEK{}~:
\[
\mu = U \approx 0{,}44~;\quad
H = -\left(2 \times \tfrac{4}{9}\log_2\tfrac{4}{9} + \tfrac{1}{9}\log_2\tfrac{1}{9}\right)
  \approx 1{,}040 + 0{,}352 \approx 1{,}39~\text{bit}.
\]

\textbf{Analyse} (confiance élevée). La banque connaît bien ses prestataires
«~historiques~», ceux de l'incident de 2026 et de l'audit de 2028. Les
quatre \NEK{} sont des prestataires que personne n'a regardés sous l'angle
de la sécurité, parce qu'ils ne sont pas «~informatiques~». Le transporteur
de fonds et l'imprimeur des chéquiers ont pourtant accès à des
informations sensibles. C'est le périmètre de complaisance du prologue,
déplacé vers l'écosystème.

\medskip
\hfill\lecture{0{,}44}{0{,}44}{1{,}39}
\end{boxlecture}

% ============================================================================
\section{Auditer un Service en Nuage}
\label{sec:c15-nuage}
% ============================================================================

\subsection{La responsabilité partagée}

Dans un service en nuage, les mesures de sécurité se répartissent entre le
fournisseur et le client, selon le modèle de service. Le fournisseur répond
de ce qu'il exploite~; le client, de ce qu'il configure et de ce qu'il y
met.

\begin{tableaucours}{L{4.0cm}ccc}{%
    \tableauHeader{Couche} &
    \tableauHeader{Infrastructure (IaaS)} &
    \tableauHeader{Plateforme (PaaS)} &
    \tableauHeader{Logiciel (SaaS)}
}
Centre de données, matériel & Fournisseur & Fournisseur & Fournisseur \\
Système d'exploitation & Client & Fournisseur & Fournisseur \\
Application & Client & Client & Fournisseur \\
Configuration du service & Client & Client & Client \\
Identités et accès des utilisateurs & Client & Client & Client \\
Données & Client & Client & Client \\
\end{tableaucours}
\vspace{-0.8em}
\begin{center}{\small\itshape Tableau~15.2. Le partage des responsabilités
(schéma général~; chaque contrat le précise)}\end{center}

Deux lignes restent toujours au client~: les identités et les données. Une
messagerie en nuage parfaitement sécurisée par son fournisseur reste
vulnérable si les comptes des agents n'ont pas de second facteur, ou si un
agent partage un dossier de clients avec tout Internet.

\subsection{Les référentiels du nuage}

L'ISO/IEC~27017 complète l'ISO/IEC~27002 par des recommandations propres
aux services en nuage, pour le fournisseur comme pour le
client~\cite{iso27017_2015}~; l'ISO/IEC~27018 traite de la protection des
données personnelles dans le nuage public~\cite{iso27018_2019}. Pour
l'auditeur, la preuve vient le plus souvent des rapports d'assurance du
fournisseur (type~2) et de ses certifications~: il vérifie qu'ils couvrent
bien le service utilisé, la période, et le centre de données qui héberge
les données.

\begin{boxscenario}{Opération SANAGA. La messagerie dans le nuage}
\textbf{Faits.} La messagerie migre vers NUAGEPRO en mai 2029. NUAGEPRO
fournit une certification ISO/IEC~27001 et un rapport de type~2.

\medskip
\textbf{Analyse} (confiance moyenne). Ces pièces couvrent la part du
fournisseur. La part de la banque se vérifie autrement~: second facteur sur
tous les comptes (une extraction de la console d'administration du
service)~; règles de partage externe désactivées~; journaux d'accès
récupérés dans le collecteur de la banque~; localisation des données et
conditions de leur transfert, à examiner au regard de la loi
n°~2024/017 (à vérifier). Le «~rapport du fournisseur~» ne remplace aucune de
ces quatre preuves.
\end{boxscenario}

\begin{boxinfo}
\textbf{La loi et les tiers.} L'article~26, alinéa~1, de la loi
n°~2010/012 vise les risques liés aux services offerts « directement ou
indirectement » (Chapitre~\ref{chap:iso27005})~\cite{loi2010012}.
Analyse (lecture du cours)~: confier un service à un prestataire ne fait
pas sortir ses risques du champ de l'obligation de l'exploitant.
\end{boxinfo}

% ============================================================================
\section{L'Intelligence Artificielle}
\label{sec:c15-ia}
% ============================================================================

\subsection{Des risques nouveaux}

Un système d'intelligence artificielle est un logiciel dont le comportement
est appris à partir de données plutôt qu'entièrement programmé. Il ajoute
aux risques classiques des risques propres~:
\begin{itemize}
    \item \textbf{la fuite de données}~: ce que l'on confie à un assistant
          public peut être conservé, réutilisé, ou exposé~;
    \item \textbf{la manipulation}~: des données d'entraînement empoisonnées,
          ou des instructions cachées dans un document soumis à un
          assistant, peuvent détourner le système~;
    \item \textbf{l'erreur opaque}~: un modèle peut se tromper sans que l'on
          sache pourquoi, et ses performances peuvent se dégrader quand les
          données changent~;
    \item \textbf{le biais}~: un modèle peut traiter différemment des
          catégories de personnes, ce qui pose des questions de droit et de
          réputation.
\end{itemize}

\subsection{Deux cadres}

\begin{tableaucours}{L{2.6cm}L{4.8cm}L{4.8cm}}{%
    \tableauHeader{} &
    \tableauHeader{ISO/IEC~42001:2023} &
    \tableauHeader{NIST AI RMF~1.0 (2023)}
}
Nature & Système de management de l'IA, certifiable~\cite{iso42001_2023} &
Cadre volontaire de gestion des risques~\cite{nist2023airmf} \\
Structure & Même structure que l'ISO/IEC~27001 (articles~4 à~10), avec une
annexe de mesures & Quatre fonctions~: gouverner, cartographier, mesurer,
gérer \\
Apport & S'intègre au SMSI existant & Décrit les caractéristiques d'une IA
digne de confiance~: valide et fiable, sûre, sécurisée et résiliente,
responsable et transparente, explicable, respectueuse de la vie privée,
équitable \\
\end{tableaucours}
\vspace{-0.8em}
\begin{center}{\small\itshape Tableau~15.3. Deux cadres pour l'IA}\end{center}

\subsection{Deux usages à la BANQUE SANAGA}

\begin{boxscenario}{Opération SANAGA. Un modèle acheté, un assistant toléré}
\textbf{Le modèle de FRAUDSCAN.} Il attribue à chaque virement un score de
risque~; au-delà d'un seuil, le virement est bloqué pour vérification.
Questions de l'auditeur~: sur quelles données le modèle a-t-il été
entraîné~? Qui a validé le seuil, et sur quels essais~? Comment la banque
sait-elle, mois après mois, que le taux de fraudes manquées ne se dégrade
pas~? Que peut-elle expliquer à un client dont le virement est bloqué, ou
au superviseur~? À ce jour, la banque ne peut répondre qu'à la deuxième
question~: les trois autres sont \NEK{}.

\medskip
\textbf{L'assistant public.} Une revue des journaux de navigation montre
que 47 agents ont soumis des textes à un assistant d'IA générative public
en un mois, dont certains contenaient des noms et des numéros de compte de
clients. Aucune règle ne l'interdisait, aucune ne l'encadrait.

\medskip
\textbf{Analyse} (confiance élevée). Le premier usage est un risque de
\emph{tiers} et de \emph{modèle}~: il relève à la fois du cycle de vie des
prestataires et d'une surveillance des performances. Le second est un
risque de \emph{fuite}, créé par un usage non déclaré~: il appelle une
règle d'usage, une solution autorisée, et une sensibilisation. Ni l'un ni
l'autre ne figurait au registre.
\end{boxscenario}

\begin{boxobjet}{R-IA-01 · Fuite de données de clients par un assistant public}
\textbf{Valeur métier~:} la confidentialité des données des clients.

\textbf{Bien support~:} navigateurs des agents, assistant public non
contractualisé.

\textbf{Événement redouté~:} divulgation de données de clients à un tiers
sans contrat~; manquement possible aux obligations de protection des
données (loi n°~2024/017, à vérifier).

\textbf{Vraisemblance~:} 4 (usage constaté, 47 agents en un mois). Base~:
journaux de navigation (\OUI{}).

\textbf{Impact~:} 3 (atteinte à la confidentialité, réputation, risque
réglementaire). Base~: appréciation du DPD (\OUI{}).

\textbf{Traitement~:} règle d'usage de l'IA approuvée par le comité de
sécurité~; blocage de l'assistant public~; mise à disposition d'un assistant
contractualisé, sans conservation des données~; sensibilisation ciblée.

\textbf{Propriétaire~:} le directeur des opérations, avec le DPD.
\end{boxobjet}

% ============================================================================
\section{L'Audit Continu}
\label{sec:c15-continu}
% ============================================================================

\subsection{De l'échantillon à la population}

Le Chapitre~\ref{chap:grc-pratique} a présenté le contrôle continu de la
\emph{conformité}~: des tests automatiques sur les contrôles. L'\termecle{audit
continu}\index{audit continu} va plus loin~: il analyse les \emph{données}
et les \emph{transactions} elles-mêmes, sur toute la population et de
manière répétée, pour y détecter les anomalies. Là où la Chambre des
comptes examinait 40 paiements de l'ANPP, l'audit continu les examine tous,
chaque mois.

\begin{boxscenario}{Opération ANPP. Tous les paiements, chaque mois}
\textbf{Le test.} Chaque mois, un programme rapproche les paiements du
fichier des agents en activité et signale tout paiement à une personne
absente de ce fichier à la date du paiement. Le programme appartient à
l'audit interne de l'ANPP, qui le fait tourner sur une copie des données.

\medskip
\textbf{Juin 2029.} Le test signale 37 exceptions. Après examen~: 29 sont
justifiées (départs enregistrés avec retard, rappels de salaire dus), 6 sont
des paiements indus, et 2 restent sans explication, les dossiers n'ayant pas
été retrouvés.
\end{boxscenario}

\begin{boxlecture}{Les 37 exceptions de juin 2029}
\textbf{Constatation.} 29 justifiées (\OUI{}), 6 indues (\NON{}), 2 sans
explication (\NEK{})~:
\[
\mu \approx 0{,}78~;\quad U \approx 0{,}05~;\quad
H \approx 0{,}276 + 0{,}426 + 0{,}228 \approx 0{,}93~\text{bit}.
\]

\textbf{Analyse} (confiance moyenne). La plupart des exceptions sont des
faux positifs, et c'est normal~: un test exhaustif signale tout ce qui
s'écarte de la règle, y compris ce qui s'explique. Les six paiements indus
sont des constats établis sur toute la population, plus solides que ceux
qu'un échantillon de 40 aurait pu fournir. Les deux \NEK{} sont l'objet
même de l'audit~: ce sont les cas où l'explication manque.

\medskip
\hfill\lecture{0{,}78}{0{,}05}{0{,}93}
\end{boxlecture}

\subsection{Les conditions et les limites}

\begin{boxprocedure}
\textbf{Mettre en place un test d'audit continu.} Le papier de travail
décrit~:
\begin{itemize}
    \item la règle testée et son critère (texte, procédure)~;
    \item les sources de données, leur exhaustivité et leur qualité (le
          fichier des agents est-il à jour~? les paiements sont-ils tous
          extraits~?)~;
    \item le programme, sa version, et sa validation (essai sur des cas
          connus)~;
    \item la fréquence, le seuil d'exception et le circuit de traitement
          des exceptions~;
    \item qui détient le programme~: l'audit interne, et non la direction
          auditée, pour préserver l'indépendance.
\end{itemize}
\end{boxprocedure}

\begin{boxalert}
\textbf{Violation typique.} Confier le test d'audit continu à la direction
qui exécute les paiements, puis lire l'absence d'exception comme une preuve.

\smallskip
\textbf{Règle~:} un test d'audit continu ne vaut que par l'exhaustivité de
ses données et l'indépendance de celui qui le fait tourner. Le même
programme, géré par la direction, est un contrôle de premier ou de deuxième
niveau (Chapitre~\ref{chap:grc}), utile, mais qui ne remplace pas l'audit.
Et trop d'exceptions non traitées finissent par ne plus être lues.
\end{boxalert}

% ============================================================================
\section{Ce que les Nouveaux Objets demandent à chaque Métier}
\label{sec:c15-metiers}
% ============================================================================

\begin{boxmetier}{Responsable GRC}
\textbf{Un registre des tiers}

Tenir la liste des prestataires critiques, leur criticité, la preuve
attendue de chacun et son statut. Y faire entrer ceux qui ne sont pas
«~informatiques~», mais touchent à l'information.
\end{boxmetier}

\begin{boxmetier}{Auditeur en cybersécurité}
\textbf{Lire un rapport d'assurance}

Vérifier que le rapport de type~2 couvre le service utilisé, la période et
le lieu des données~; repérer les mesures que le fournisseur attend du
client (souvent listées dans le rapport) et les tester chez le client.
Pour l'IA, demander les preuves d'entraînement, de validation et de
surveillance du modèle.
\end{boxmetier}

\begin{boxmetier}{RSSI}
\textbf{Encadrer les usages avant de les interdire}

Un usage non déclaré de l'IA révèle un besoin. La RSSI propose une règle
d'usage, une solution autorisée et une sensibilisation~; le blocage seul
déplace l'usage ailleurs.
\end{boxmetier}

\begin{boxmetier}{Gestionnaire des risques cyber}
\textbf{Faire entrer l'IA dans le registre}

Chaque système d'IA en usage est inventorié, avec son propriétaire, ses
données, son fournisseur et ses scénarios de risque~: fuite, manipulation,
erreur, biais.
\end{boxmetier}

\separateur

\begin{boxresume}
\textbf{Ce qu'il faut retenir de ce chapitre~:}
\begin{itemize}
    \item \textbf{Les tiers} sont une porte d'entrée, un risque de
          défaillance et une chaîne~; le risque reste celui de
          l'organisation.
    \item \textbf{Le cycle de vie d'un tiers}~: sélection, contrat,
          exploitation, sortie, avec une preuve à chaque étape.
    \item \textbf{Dans le nuage}, les identités et les données restent
          toujours au client~; le rapport du fournisseur ne couvre que sa
          part.
    \item \textbf{L'IA} ajoute fuite, manipulation, erreur opaque et biais~;
          ISO/IEC~42001 et le NIST AI RMF l'encadrent.
    \item \textbf{L'audit continu} examine toute la population~; il vaut par
          l'exhaustivité des données et l'indépendance de son exploitant.
\end{itemize}
\end{boxresume}

\begin{boxinfo}
\textbf{Points de contrôle.} Vous avez compris ce chapitre si vous savez~:
\begin{itemize}
    \item répartir une mesure entre fournisseur et client pour un service
          donné~;
    \item dire ce qu'un rapport d'assurance prouve et ce qu'il laisse au
          client~;
    \item apprécier un usage de l'IA sous forme de fiche de risque~;
    \item concevoir un test d'audit continu et ses conditions de validité.
\end{itemize}
\end{boxinfo}

% ============================================================================
\section*{Exercices}
% ============================================================================

\begin{boxexo}[title={\ding{45}\quad Exercice~15.1. Qui en répond~? \niveaubase}]
La messagerie de la BANQUE SANAGA est un service en nuage de type logiciel
(SaaS). Pour chacune des tâches suivantes, dites si elle revient à NUAGEPRO
ou à la banque.
\begin{enumerate}[(a)]
    \item Appliquer les correctifs du système d'exploitation des serveurs.
    \item Activer le second facteur pour les agents.
    \item Désactiver le partage externe des documents.
    \item Protéger physiquement le centre de données.
    \item Désactiver le compte d'un agent parti.
    \item Conserver les journaux d'accès selon la loi.
\end{enumerate}
\end{boxexo}

\begin{boxexo}[title={\ding{45}\quad Exercice~15.2. Le contrat de FRAUDSCAN \niveaumoyen}]
\begin{enumerate}[(a)]
    \item Rédigez cinq clauses de sécurité pour le contrat du modèle de
          détection de fraude, en couvrant les quatre étapes du cycle de vie
          d'un tiers.
    \item Ajoutez deux clauses propres à l'IA (données d'entraînement,
          surveillance des performances).
    \item Quelles preuves la banque devra-t-elle obtenir chaque année pour
          lever les trois \NEK{} du modèle~?
\end{enumerate}
\end{boxexo}

\begin{boxexo}[title={\ding{45}\quad Exercice~15.3. L'IA chez LiZbiZ \niveaumoyen}]
LiZbiZ envisage un modèle qui optimise les tournées de livraison à partir
des données de suivi des expéditions et des adresses des clients.
\begin{enumerate}[(a)]
    \item Identifiez deux scénarios de risque propres à l'IA, et rédigez la
          fiche de l'un d'eux.
    \item Placez les actions de LiZbiZ dans les quatre fonctions du NIST AI
          RMF.
    \item PHARMACAM demande si ce modèle « est conforme à
          l'ISO/IEC~42001 ». Que répondez-vous~?
\end{enumerate}
\end{boxexo}

\begin{boxexo}[title={\ding{45}\quad Exercice~15.4. Un audit continu des ajustements de stock \niveauexpert}]
Le cours ACSI a établi que 148~millions FCFA d'ajustements de stock, dont
71~\% en fin de mois, avaient été enregistrés sous le compte de M.~FOUDA à
Douala.
\begin{enumerate}[(a)]
    \item Concevez un test d'audit continu qui aurait signalé ces
          ajustements~: règle, données, seuils, fréquence.
    \item Rédigez le papier de travail de mise en place (exhaustivité,
          validation, indépendance).
    \item Lors du premier mois, le test signale 120 exceptions. Comment
          éviter que les exceptions ne soient plus lues, sans perdre les
          vraies anomalies~?
    \item Un mois, 80 exceptions sont justifiées, 15 confirmées comme
          anormales, 5 sans explication. Calculez $\mu$, $U$ et $H$, et dites
          ce que chaque catégorie appelle.
\end{enumerate}
\end{boxexo}

\bigskip

\begin{boxinfo}
\textbf{Pour aller plus loin.}
\begin{itemize}
    \item L'ISO/IEC~27002:2022~\cite{iso27002_2022}, mesures 5.19 à 5.23.
    \item L'ISO/IEC~27017~\cite{iso27017_2015} et l'ISO/IEC~27018~\cite{iso27018_2019}.
    \item L'ISO/IEC~42001:2023~\cite{iso42001_2023} et le NIST AI
          RMF~\cite{nist2023airmf}.
\end{itemize}
\end{boxinfo}

\begin{boxinfo}
\textbf{Transition.} Tiers opaques, modèles qui apprennent, populations
entières analysées~: à chaque fois, la question de l'auditeur est la même~:
que sait-on, et que ne sait-on pas~? La Partie~VI, consacrée à la
recherche, montre comment l'auditologie fait de cette question un objet de
mesure. C'est le Chapitre~\ref{chap:auditologie}.
\end{boxinfo}
